\documentclass[11pt,oneside]{article}

\usepackage[T1]{fontenc}
\usepackage[utf8]{inputenc}
\usepackage{lmodern}
\usepackage{microtype}
\usepackage[a4paper,margin=28mm,headheight=15pt]{geometry}
\usepackage{amsmath,amssymb,amsthm,mathtools}
\usepackage{booktabs,longtable,array,tabularx}
\usepackage{enumitem}
\usepackage{xcolor}
\usepackage{hyperref}
\usepackage{cleveref}
\usepackage{fancyhdr}
\usepackage{listings}
\usepackage{graphicx}
\usepackage{tikz}
\usetikzlibrary{arrows.meta,positioning,calc,fit}

\definecolor{linkblue}{RGB}{28,73,120}
\definecolor{codegray}{RGB}{245,245,245}
\hypersetup{
  colorlinks=true,
  linkcolor=linkblue,
  citecolor=linkblue,
  urlcolor=linkblue,
  pdftitle={Additive progress on the Alaoglu--Erdos conjecture},
  pdfauthor={Research working note prepared for Vladimir Reshetnikov}
}

\pagestyle{fancy}
\fancyhf{}
\lhead{Alaoglu--Erd\H{o}s: additive progress report B}
\rhead{\thepage}

\setlist[itemize]{topsep=4pt,itemsep=2pt,parsep=1pt}
\setlist[enumerate]{topsep=4pt,itemsep=2pt,parsep=1pt}
\setlength{\parindent}{1.25em}
\setlength{\parskip}{0.25em}
\allowdisplaybreaks

\lstset{
  basicstyle=\ttfamily\small,
  backgroundcolor=\color{codegray},
  frame=single,
  breaklines=true,
  columns=fullflexible,
  keepspaces=true,
  showstringspaces=false
}

\newtheorem{theorem}{Theorem}[section]
\newtheorem{proposition}[theorem]{Proposition}
\newtheorem{lemma}[theorem]{Lemma}
\newtheorem{corollary}[theorem]{Corollary}
\newtheorem{conjecture}[theorem]{Conjecture}
\newtheorem{target}[theorem]{Target lemma}
\newtheorem{warning}[theorem]{Proof-audit warning}
\theoremstyle{definition}
\newtheorem{definition}[theorem]{Definition}
\newtheorem{experiment}[theorem]{Computational experiment}
\theoremstyle{remark}
\newtheorem{remark}[theorem]{Remark}

\newcommand{\Q}{\mathbb{Q}}
\newcommand{\R}{\mathbb{R}}
\newcommand{\Z}{\mathbb{Z}}
\newcommand{\F}{\mathbb{F}}
\newcommand{\Gm}{\mathbb{G}_{\mathrm m}}
\newcommand{\ord}{\operatorname{ord}}
\newcommand{\rank}{\operatorname{rank}}
\newcommand{\Gal}{\operatorname{Gal}}
\newcommand{\Frob}{\operatorname{Frob}}
\newcommand{\Span}{\operatorname{span}}
\newcommand{\dens}{\operatorname{dens}}
\newcommand{\mult}{\operatorname{mult}}
\newcommand{\rad}{\operatorname{rad}}
\newcommand{\AE}{\mathrm{AE}}
\newcommand{\ellogtwo}{\log 2}
\newcommand{\ellogthree}{\log 3}

\title{\textbf{Additive Progress on the Alaoglu--Erd\H{o}s Conjecture}\\[0.4em]
\large Kummer--Chebotarev incompatibility, mantissa-orbit entropy,\\
and a rigidity barrier for auxiliary polynomials}
\author{Research working note prepared for Vladimir Reshetnikov}
\date{22 August 2026}

\begin{document}
\maketitle

\begin{abstract}
Let $x\in\R$ and suppose that $m=2^x$ and $n=3^x$ are integers.  The
Alaoglu--Erd\H{o}s conjecture asserts that $x$ must be an integer.  This
companion report is deliberately additive: it does not repeat the large
survey and proof search in Report~A.  It develops three pieces of new
infrastructure.

First, it associates to $(2,3,m,n)$ an exact order-$\ell$ Kummer
compatibility test.  Kummer theory and Chebotarev then give closed formulas
for the density of primes at which a common residue exponent exists.  For a
hypothetical nonintegral counterexample that density is asymptotic to
$1/\ell$, while incompatibility has density tending to one.  The rank-three
case, which initially appears exceptional, is shown to be transverse; its
exact compatibility density is
$\ell^{-1}-\ell^{-2}+\ell^{-3}$.

Second, the normalized rational points
\[
 U_j=2^{\{jx\}}=\frac{m^j}{2^{\lfloor jx\rfloor}},\qquad
 V_j=3^{\{jx\}}=\frac{n^j}{3^{\lfloor jx\rfloor}}
\]
are studied simultaneously in the real and finite topologies.  Conditional
on any real subarc of the exponential curve, their order-$\ell$ residue
symbols are uniformly distributed on the image of an explicit $2\times2$
matrix.  At the Chebotarev-generic primes that matrix is invertible, so the
finite residue data has full two-dimensional entropy even on an arbitrarily
short real arc.  This is a rigorous archimedean-to-finite-place theorem, but
it points in the opposite direction from a naive local-exponent proof.

Third, an exact local zero estimate is proved.  If an algebraic-coefficient
polynomial is restricted to $z\mapsto(2^z,3^z)$ at an algebraic point of the
curve, its order of contact is exactly its ordinary bivariate multiplicity.
Thus the roughly $D^2$ coefficients of a degree-$D$ auxiliary polynomial
cannot buy order-$D^2$ directional vanishing; only order $O(D)$ is possible.
This identifies, rather than merely gestures at, the obstruction in a broad
class of interpolation-determinant attacks.

The conjecture is not claimed solved here.  The report ends with three
precise bridge lemmas which would finish it, a proof-status ledger, a
finite-field verification program, and a staged Lean formalization plan.
\end{abstract}

\tableofcontents
\newpage

\section{Scope, novelty boundary, and proof status}

\subsection{The statement}

We write the conjecture in the following form.

\begin{conjecture}[Alaoglu--Erd\H{o}s]\label{conj:AE}
If $x\in\R$ and $2^x,3^x\in\Z$, then $x\in\Z$.
\end{conjecture}

Only positive integers can occur: if $x<0$, then $0<2^x<1$, while if
$x=0$ the conclusion is immediate.  Hence every putative counterexample has
$x>0$ and $m=2^x,n=3^x\in\Z_{>0}$.

The external reports about recent machine-assisted mathematical discoveries
are treated as motivation, not as evidence.  Every assertion used below is
proved in the text or reduced to a named, standard theorem.  In particular,
no unpublished announcement is used as a lemma.

\subsection{How this document avoids duplicating Report A}

The attached source corpus was decompressed and used as an exclusion map.
Its machine-audited metadata is recorded here so that the relationship
between the documents is reproducible:
\begin{center}
\begin{tabular}{@{}ll@{}}
\toprule
Source file & \texttt{alaoglu-erdos-unified-report-A.tex.gz}\\
Decompressed lines & \texttt{20,925}\\
Whitespace-delimited words & \texttt{113,833}\\
Structural headings & \texttt{491}\\
SHA-256 of compressed source & \texttt{1e77491ea5cc0911d56e}\\
\bottomrule
\end{tabular}
\end{center}

The following exact strings were also checked against Report~A.  This is only
a lexical guard against accidental repetition, not a claim of priority in
the literature.
\begin{center}
\begin{tabularx}{\textwidth}{@{}Xc@{}}
\toprule
String identifying a result developed below & Present verbatim in Report A?\\
\midrule
\texttt{mantissa-orbit entropy} & no\\
\texttt{transcendental-direction rigidity} & no\\
\texttt{fixed-t Kummer density} & no\\
\texttt{q\string^3-q+1} & no\\
\texttt{rank-three transversality} & no\\
\bottomrule
\end{tabularx}
\end{center}

The baseline reductions in \cref{sec:baseline} are included only to make the
new arguments self-contained.  The substantive additions begin with the
rank-three transversality lemma, the exact Kummer counts, the conditional
mantissa equidistribution theorem, and the auxiliary-polynomial rigidity
theorem.

\subsection{What is proved, and what is not}

The main unconditional conclusions of this report are as follows.

\begin{enumerate}
\item Every nonintegral counterexample has multiplicative rank at least three:
      $\rank\langle2,3,m,n\rangle\ge3$.
\item In rank three, the unique multiplicative relation is necessarily
      transverse between the input pair $(2,3)$ and output pair $(m,n)$.
\item For every sufficiently large odd prime $\ell$, the relative
      Chebotarev density of primes admitting a common order-$\ell$ residue
      exponent is given by an exact formula.  It is
      $\ell^{-1}-\ell^{-2}+\ell^{-3}$ in rank three and
      $\ell^{-1}-\ell^{-3}+\ell^{-4}$ in rank four.
\item At a density $1-O(1/\ell)$ of primes, the normalized mantissa orbit has
      all $\ell^2$ residue-symbol pairs, uniformly, even after conditioning
      on an arbitrarily short real interval.
\item Algebraic auxiliary polynomials have no anomalously high directional
      contact with the curve $z\mapsto(2^z,3^z)$ at its algebraic points.
\item A common prime divisor of the simultaneous near-power differences
      $m^b-2^a$ and $n^b-3^a$ is forced into the rare Kummer-compatible
      Chebotarev set whenever $\ell\nmid b$.
\end{enumerate}

None of these statements by itself proves \cref{conj:AE}.  The precise
remaining gaps are isolated in \cref{sec:bridges}.  The most important
negative conclusion is also constructive: simply transporting the real
exponent $x$ to a $p$-adic or finite-field exponent cannot work at generic
primes, because a hypothetical counterexample is provably incompatible at
almost all of them.

\section{Baseline reductions and the multiplicative-rank floor}\label{sec:baseline}

Throughout the rest of the report, a \emph{counterexample datum} means a
triple
\[
 (x,m,n),\qquad x\in\R\setminus\Z,\qquad m=2^x\in\Z_{>0},\quad
 n=3^x\in\Z_{>0}.
\]
The point of introducing the term is not to assume that such a datum exists,
but to make conditional deductions compact.

\subsection{Rational and algebraic exponents are already excluded}

\begin{proposition}\label{prop:rational-integer}
If $x\in\Q$ and $2^x\in\Z$, then $x\in\Z$.
\end{proposition}

\begin{proof}
Write $x=a/b$ with $a\in\Z$, $b\in\Z_{>0}$, and $\gcd(a,b)=1$.  Since
$2^{a/b}$ is a positive integer $m$, one has $2^a=m^b$.  Comparing the
$2$-adic valuations gives $a=bv_2(m)$.  Hence $b\mid a$, so $b=1$.
\end{proof}

\begin{corollary}\label{cor:x-transcendental}
In a counterexample datum, $x$ is transcendental.
\end{corollary}

\begin{proof}
By \cref{prop:rational-integer}, $x$ is irrational.  If it were algebraic,
then the Gelfond--Schneider theorem, applied to the algebraic base $2$ and
the algebraic irrational exponent $x$, would make every value of $2^x$
transcendental.  This contradicts $2^x=m\in\Z$.
\end{proof}

Set
\[
 \alpha=\frac{\log3}{\log2}.
\]

\begin{proposition}\label{prop:alpha-transcendental}
The number $\alpha$ is transcendental.
\end{proposition}

\begin{proof}
It is not rational, since $\alpha=r/s\in\Q$ would imply $2^r=3^s$ after
clearing denominators.  If $\alpha$ were algebraic irrational, then
Gelfond--Schneider would make $2^\alpha$ transcendental, whereas
$2^\alpha=3$.
\end{proof}

The conjecture can therefore be viewed as a very special rank-one relation
between a $2\times2$ matrix of logarithms of positive integers:
\begin{equation}\label{eq:logdet}
 \det\begin{pmatrix}
  \log2&\log3\\
  \log m&\log n
 \end{pmatrix}=0,
 \qquad
 (\log m,\log n)=x(\log2,\log3).
\end{equation}
The Four Exponentials Conjecture would rule out this matrix when $x$ is
irrational.  The integrality of $m,n$ is the extra datum one hopes to exploit.

\subsection{Multiplicative rank}

For a finitely generated subgroup $\Gamma\subset\Q_{>0}^{\times}$, its rank
means its rank as a free abelian group.  Put
\[
 \Gamma=\langle2,3,m,n\rangle\subset\Q_{>0}^{\times},\qquad r=\rank\Gamma.
\]
Since $2$ and $3$ are multiplicatively independent, $r\ge2$.

\begin{proposition}[rank floor]\label{prop:rank-floor}
For a counterexample datum one has $r\ge3$.
\end{proposition}

\begin{proof}
Suppose $r=2$.  Since the valuation vectors of $2$ and $3$ are independent,
any prime outside $\{2,3\}$ occurring in $m$ or $n$ would raise the rank.
Consequently
\[
 m=2^a3^b,\qquad n=2^c3^d
\]
with $a,b,c,d\in\Z_{\ge0}$.  Dividing logarithms by $\log2$ gives
\[
 x=a+b\alpha,\qquad x\alpha=c+d\alpha.
\]
Eliminating $x$ yields
\[
 b\alpha^2+(a-d)\alpha-c=0.
\]
By \cref{prop:alpha-transcendental}, all three rational coefficients vanish.
Thus $b=c=0$ and $a=d$, whence $x=a\in\Z$, contrary to the definition of a
counterexample datum.
\end{proof}

\begin{corollary}\label{cor:outside-prime}
In a counterexample datum, at least one of $m,n$ has a prime divisor outside
$\{2,3\}$.
\end{corollary}

\begin{proof}
If both were $\{2,3\}$-smooth, then $\Gamma$ would have rank two, contrary to
\cref{prop:rank-floor}.
\end{proof}

\subsection{The rank-three normal form and its transversality}

If $r=3$, the homomorphism $\Z^4\to\Gamma$ sending the standard basis to
$2,3,m,n$ has a rank-one kernel.  Let
\begin{equation}\label{eq:primitive-relation}
 2^A3^B m^C n^D=1,
 \qquad (A,B,C,D)\in\Z^4
\end{equation}
be a primitive generator of that kernel, unique up to sign.
Substitution of $m=2^x,n=3^x$ gives
\begin{equation}\label{eq:mobius-relation}
 (A+Cx)+(B+Dx)\alpha=0,
 \qquad
 x=-\frac{A+B\alpha}{C+D\alpha}.
\end{equation}
The denominator is nonzero: otherwise $C+D\alpha=0$, and the transcendence
of $\alpha$ would force $C=D=0$, which would make
$2^A3^B=1$ and hence $A=B=0$.

\begin{lemma}[rank-three transversality]\label{lem:rank3-transverse}
For a rank-three counterexample datum, the primitive relation
\eqref{eq:primitive-relation} satisfies
\[
 AD-BC\ne0.
\]
Equivalently, the coefficient vectors $(A,B)$ and $(C,D)$ are not
proportional over $\Q$.
\end{lemma}

\begin{proof}
If $(A,B)=\lambda(C,D)$ for some $\lambda\in\Q$, then
\eqref{eq:mobius-relation} becomes
\[
 (x+\lambda)(C+D\alpha)=0.
\]
The second factor is nonzero, so $x=-\lambda\in\Q$.  By
\cref{prop:rational-integer}, $x\in\Z$, a contradiction.
\end{proof}

This tiny determinant will be responsible for the clean rank-three
Chebotarev formula.  It removes an apparent exceptional case that would
otherwise double the compatible density.

\section{Kummer residue coordinates}\label{sec:kummer-setup}

\subsection{The Kummer space}

Fix an odd prime $\ell$ and put
\[
 K_\ell=\Q(\zeta_\ell),
\]
where $\zeta_\ell$ is a chosen primitive $\ell$th root of unity.  Let
$V_\ell$ be the $\F_\ell$-subspace of
$K_\ell^\times/K_\ell^{\times\ell}$ generated by the classes of
$2,3,m,n$, and write
\[
 r_\ell=\dim_{\F_\ell}V_\ell.
\]
Kummer theory identifies
\begin{equation}\label{eq:kummer-galois}
 \Gal\bigl(K_\ell(2^{1/\ell},3^{1/\ell},m^{1/\ell},n^{1/\ell})/K_\ell\bigr)
 \cong \operatorname{Hom}_{\F_\ell}(V_\ell,\mu_\ell),
\end{equation}
so this Galois group has order $\ell^{r_\ell}$.

The dimension $r_\ell$ agrees with the ordinary multiplicative rank for all
but finitely many $\ell$.  The finite exceptional set is worth stating
correctly; ignoring it is a common source of overstrong Kummer claims.

\begin{lemma}[eventual Kummer rank]\label{lem:eventual-rank}
Let $\Gamma=\langle2,3,m,n\rangle$ have rank $r$.  There is a positive
integer $I_\Gamma$ such that, for every odd prime $\ell\nmid I_\Gamma$,
\[
 r_\ell=r.
\]
\end{lemma}

\begin{proof}
Let $S$ be the finite set of rational primes occurring in $6mn$.  The
valuation map identifies the group of positive rationals supported on $S$
with the lattice $\Z^S$.  The image $L$ of $\Gamma$ is a rank-$r$ sublattice.
Let
\[
 L^{\mathrm{sat}}=(L\otimes_\Z\Q)\cap\Z^S,
 \qquad I_\Gamma=[L^{\mathrm{sat}}:L].
\]
For $\ell\nmid I_\Gamma$, reduction modulo $\ell$ gives an $r$-dimensional
image in $\Q^\times/\Q^{\times\ell}$.

It remains to see that passing from $\Q$ to $K_\ell$ creates no new
$\ell$th-power relation among positive rationals.  Suppose
$a\in\Q_{>0}$ and $a=y^\ell$ for some $y\in K_\ell$.  At a rational prime
$p\ne\ell$, the cyclotomic extension is unramified, so every valuation of
$a$ above $p$ is $v_p(a)$ and must be divisible by $\ell$.  At $p=\ell$ the
ramification index is $\ell-1$; hence $\ell\mid(\ell-1)v_\ell(a)$, and
coprimality gives $\ell\mid v_\ell(a)$.  All rational prime valuations of
$a$ are therefore divisible by $\ell$, and positivity gives
$a\in\Q^{\times\ell}$.  Thus the rational Kummer classes inject into the
cyclotomic Kummer space, proving the claim.
\end{proof}

We call an odd prime $\ell$ \emph{good for the datum} if it avoids
$I_\Gamma$ and, in rank three, also avoids the nonzero integer $AD-BC$ from
\cref{lem:rank3-transverse}.  Only finitely many odd primes are not good.

\subsection{Residue-symbol coordinates}

Let
\[
 L_\ell=K_\ell(2^{1/\ell},3^{1/\ell},m^{1/\ell},n^{1/\ell}).
\]
For a prime ideal $\mathfrak p$ of $K_\ell$ unramified in $L_\ell$ and not
dividing $6mn\ell$, define $\kappa_{\mathfrak p}(a)\in\F_\ell$ by
\[
 \left(\frac{a}{\mathfrak p}\right)_\ell
   =\zeta_\ell^{\kappa_{\mathfrak p}(a)},
 \qquad a\in\{2,3,m,n\}.
\]
Equivalently, $\kappa_{\mathfrak p}$ is the Kummer character represented by
$\Frob_{\mathfrak p}$.  Put
\begin{equation}\label{eq:kappa-vector}
 \boldsymbol\kappa_{\mathfrak p}
 =\bigl(u_1,u_2,v_1,v_2\bigr)
 =\bigl(\kappa_{\mathfrak p}(2),\kappa_{\mathfrak p}(3),
        \kappa_{\mathfrak p}(m),\kappa_{\mathfrak p}(n)\bigr)
 \in\F_\ell^4.
\end{equation}
Changing the chosen primitive root $\zeta_\ell$ multiplies all four
coordinates by one nonzero scalar, so all zero/nonzero and proportionality
conditions used below are canonical.

By Chebotarev, the vectors \eqref{eq:kappa-vector} are uniformly distributed
in the $r_\ell$-dimensional linear subspace
\[
 W_\ell=V_\ell^\vee\subset\F_\ell^4.
\]
More explicitly, every vector in $W_\ell$ occurs on a set of prime ideals of
$K_\ell$ of relative density $\ell^{-r_\ell}$.

\subsection{Exact compatibility versus a vanishing determinant}

There are two related local tests, and they must not be conflated.

\begin{definition}[exact common residue exponent]\label{def:exact-compatible}
A vector $(u_1,u_2,v_1,v_2)\in\F_\ell^4$ is \emph{compatible} if there exists
$t\in\F_\ell$ such that
\[
 (v_1,v_2)=t(u_1,u_2).
\]
Let $\mathcal C_\ell$ denote the set of compatible vectors, and let
\[
 \mathcal C_{\ell,t}=\{(u_1,u_2,tu_1,tu_2):u_1,u_2\in\F_\ell\}
\]
be the fixed-$t$ compatibility plane.
\end{definition}

The projective determinant
\begin{equation}\label{eq:Delta}
 \Delta(u_1,u_2,v_1,v_2)=u_1v_2-u_2v_1
\end{equation}
vanishes on $\mathcal C_\ell$, but the converse fails when
$(u_1,u_2)=(0,0)$ and $(v_1,v_2)\ne(0,0)$.  Therefore
$\Delta=0$ is a useful rank detector, not an exact common-exponent detector.
The distinction changes the main term beyond first order and is important in
formal verification.

\begin{lemma}[finite compatibility count]\label{lem:C-count}
One has
\[
 |\mathcal C_\ell|=\ell^3-\ell+1.
\]
\end{lemma}

\begin{proof}
For a nonzero source vector $u=(u_1,u_2)$, the $\ell$ values of $t$ give
$\ell$ distinct targets $tu$.  This contributes
$\ell(\ell^2-1)$ vectors.  For $u=0$, exact compatibility permits only the
zero target, contributing one more vector.
\end{proof}

\begin{warning}\label{warn:det-not-exact}
Replacing $|\mathcal C_\ell|$ by the number of singular $2\times2$ matrices,
$\ell^3+\ell^2-\ell$, silently counts the $\ell^2-1$ false positives with
zero source and nonzero target.  Any Lean formalization should define
compatibility by an existential scalar before proving a determinant
corollary.
\end{warning}

\section{Exact Kummer--Chebotarev density laws}\label{sec:kummer-density}

\subsection{The universal density formula}

The Kummer space gives a formula before any rank classification.

\begin{theorem}[Kummer compatibility density]\label{thm:universal-density}
For every odd prime $\ell$, the relative density among unramified prime
ideals $\mathfrak p$ of $K_\ell$ at which a common order-$\ell$ residue
exponent exists is
\begin{equation}\label{eq:universal-density}
 \dens_\ell(\mathrm{compatible})
 =\frac{|W_\ell\cap\mathcal C_\ell|}{\ell^{r_\ell}}.
\end{equation}
For a fixed $t\in\F_\ell$, the corresponding density is
\[
 \dens_\ell(t)=
 \frac{|W_\ell\cap\mathcal C_{\ell,t}|}{\ell^{r_\ell}}.
\]
\end{theorem}

\begin{proof}
The Kummer isomorphism \eqref{eq:kummer-galois} identifies the Galois group
with the finite additive space $W_\ell$.  The Frobenius element at
$\mathfrak p$ has coordinate vector
$\boldsymbol\kappa_{\mathfrak p}$.  Chebotarev assigns density
$1/\ell^{r_\ell}$ to each Galois element.  Summing over the indicated subsets
proves both formulas.
\end{proof}

The content is therefore finite linear algebra.  The rank-three
transversality lemma makes that linear algebra particularly clean.

\subsection{Rank four}

\begin{theorem}[rank-four law]\label{thm:rank4-law}
Suppose $r_\ell=4$.  Then, for every $t\in\F_\ell$,
\[
 \dens_\ell(t)=\frac1{\ell^2},
\]
and the exact compatibility density is
\begin{equation}\label{eq:rank4-compatible}
 \dens_\ell(\mathrm{compatible})
 =\frac{\ell^3-\ell+1}{\ell^4}
 =\frac1\ell-\frac1{\ell^3}+\frac1{\ell^4}.
\end{equation}
The incompatibility density is therefore
\begin{equation}\label{eq:rank4-incompatible}
 1-\frac1\ell+\frac1{\ell^3}-\frac1{\ell^4}.
\end{equation}
\end{theorem}

\begin{proof}
Here $W_\ell=\F_\ell^4$.  Each fixed-$t$ plane has $\ell^2$ points, giving
density $\ell^2/\ell^4=\ell^{-2}$.  The union count is
\cref{lem:C-count}; division by $\ell^4$ proves
\eqref{eq:rank4-compatible}.
\end{proof}

The formula can also be seen as a corrected union of $\ell$ planes.  Any two
planes $\mathcal C_{\ell,t}$ and $\mathcal C_{\ell,s}$ with $t\ne s$ meet
only at the zero vector.  Thus
\[
 |\mathcal C_\ell|=1+\ell(\ell^2-1).
\]
The common zero Frobenius is the only overlap.

\subsection{Rank three: the general hyperplane count}

Before using transversality, it is useful to record the exact count for an
arbitrary hyperplane.  Let
\begin{equation}\label{eq:hyperplane}
 W=\{(u_1,u_2,v_1,v_2):Au_1+Bu_2+Cv_1+Dv_2=0\}
\end{equation}
be a hyperplane in $\F_\ell^4$.  Define linear forms on
$u=(u_1,u_2)$ by
\[
 L(u)=Au_1+Bu_2,\qquad M(u)=Cu_1+Du_2.
\]
On the fixed-$t$ compatibility plane $v=tu$, the hyperplane equation becomes
\begin{equation}\label{eq:L+tM}
 (L+tM)(u)=0.
\end{equation}

\begin{proposition}[exact hyperplane count]\label{prop:hyperplane-count}
For the hyperplane \eqref{eq:hyperplane},
\[
 |W\cap\mathcal C_\ell|=
 \begin{cases}
  2\ell^2-2\ell+1,
     & M\ne0\text{ and }L\in\F_\ell M,\\[2mm]
  \ell^2-\ell+1,
     & \text{otherwise}.
 \end{cases}
\]
\end{proposition}

\begin{proof}
The zero source contributes the single zero vector.  For a nonzero source
$u$, the number of scalars $t$ solving $L(u)+tM(u)=0$ is one if
$M(u)\ne0$, is $\ell$ if $L(u)=M(u)=0$, and is zero if
$M(u)=0\ne L(u)$.

If $M\ne0$ and $L$ is proportional to $M$, their kernels coincide and contain
$\ell-1$ nonzero vectors.  There are $\ell^2-\ell$ nonzero vectors outside
that kernel, so the total is
\[
 1+(\ell^2-\ell)+\ell(\ell-1)=2\ell^2-2\ell+1.
\]
If $M\ne0$ and $L,M$ are independent, the common kernel is zero and
$M$ has $\ell-1$ nonzero zeros.  Hence the total is
$1+(\ell^2-1)-(\ell-1)=\ell^2-\ell+1$.  Finally, if $M=0$, then $L\ne0$;
only the $\ell-1$ nonzero vectors in $\ker L$ work, each with all $\ell$
values of $t$, again giving $1+\ell(\ell-1)$.
\end{proof}

The larger first case is exactly the nontransverse case.  It cannot persist
for a nonintegral rank-three counterexample.

\subsection{Rank three for a hypothetical counterexample}

\begin{theorem}[transverse rank-three law]\label{thm:rank3-law}
Let $(x,m,n)$ be a rank-three counterexample datum, and let $\ell$ be good for
the datum.  Then, for every $t\in\F_\ell$,
\[
 \dens_\ell(t)=\frac1{\ell^2},
\]
and the exact compatibility density is
\begin{equation}\label{eq:rank3-compatible}
 \dens_\ell(\mathrm{compatible})
 =\frac{\ell^2-\ell+1}{\ell^3}
 =\frac1\ell-\frac1{\ell^2}+\frac1{\ell^3}.
\end{equation}
Thus incompatibility has density
\begin{equation}\label{eq:rank3-incompatible}
 1-\frac1\ell+\frac1{\ell^2}-\frac1{\ell^3}.
\end{equation}
\end{theorem}

\begin{proof}
For a good $\ell$, $r_\ell=3$ and $W_\ell$ is the reduction modulo $\ell$ of
the primitive relation \eqref{eq:primitive-relation}:
\[
 Au_1+Bu_2+Cv_1+Dv_2=0.
\]
Goodness includes $\ell\nmid AD-BC$, so $L$ and $M$ are independent.  For
each fixed $t$, the form $L+tM$ is nonzero and has a one-dimensional kernel,
so $|W_\ell\cap\mathcal C_{\ell,t}|=\ell$.  Since
$|W_\ell|=\ell^3$, the fixed-$t$ density is $\ell^{-2}$.

For distinct $s,t$, the sets
$W_\ell\cap\mathcal C_{\ell,s}$ and
$W_\ell\cap\mathcal C_{\ell,t}$ meet only at zero.  Their union therefore
has
\[
 1+\ell(\ell-1)=\ell^2-\ell+1
\]
points, proving \eqref{eq:rank3-compatible}.
\end{proof}

\begin{remark}[fixed exponents are uniformly expensive]\label{rem:fixed-t}
The fixed-$t$ density is exactly $\ell^{-2}$ in both ranks three and four.
Rank changes only the density of the completely split zero vector and hence
the overlap correction when the $\ell$ possible exponents are united.
This makes the main term $1/\ell$ transparent: there are $\ell$ candidate
exponents, each imposing two independent residue equations.
\end{remark}

\subsection{The projective determinant law}

The determinant $\Delta$ from \eqref{eq:Delta} is not the exact detector, but
it will control the finite entropy of the mantissa orbit.  Its density is
also elementary.

\begin{proposition}[projective rank law]\label{prop:projective-law}
Let $\ell$ be an odd prime.
\begin{enumerate}
\item If $r_\ell=4$, then
\[
 \dens_\ell(\Delta=0)
 =\frac{\ell^3+\ell^2-\ell}{\ell^4}
 =\frac1\ell+\frac1{\ell^2}-\frac1{\ell^3}.
\]
\item If $(x,m,n)$ is a rank-three counterexample datum and $\ell$ is good,
then
\[
 \dens_\ell(\Delta=0)=\frac1\ell.
\]
\end{enumerate}
\end{proposition}

\begin{proof}
In rank four, $\Delta=0$ means that a $2\times2$ matrix has rank at most one.
There is one zero matrix and
\[
 \frac{(\ell^2-1)^2}{\ell-1}
\]
rank-one matrices, for a total of $\ell^3+\ell^2-\ell$.

In rank three, $W_\ell$ is the hyperplane whose coefficient matrix
$\bigl(\begin{smallmatrix}A&B\\C&D\end{smallmatrix}\bigr)$ is invertible
modulo $\ell$.  Left and right changes of basis preserve singularity and
reduce the hyperplane to $a+d=0$.  A matrix
$\bigl(\begin{smallmatrix}a&b\\c&-a\end{smallmatrix}\bigr)$ is singular when
$a^2+bc=0$.  If $b\ne0$, there are $\ell(\ell-1)$ choices, and if $b=0$ one
must have $a=0$ while $c$ is arbitrary, giving another $\ell$ choices.  The
total is $\ell^2$ out of $\ell^3$.
\end{proof}

\begin{corollary}[generic local incompatibility]\label{cor:generic-incompatibility}
For every hypothetical counterexample datum and every sufficiently large odd
prime $\ell$, exact common-exponent compatibility occurs on a relative
Chebotarev set of density $1/\ell+O(1/\ell^2)$, while projective rank collapse
occurs on a set of density $1/\ell+O(1/\ell^2)$.  In particular, both kinds
of local alignment fail on a density tending to one as $\ell\to\infty$.
\end{corollary}

\begin{proof}
By \cref{prop:rank-floor}, the rank is three or four.  Apply
\cref{thm:rank4-law,thm:rank3-law,prop:projective-law} outside the finite set
of bad primes.
\end{proof}

\subsection{Why this is an obstruction, not yet a contradiction}

The real equation $(\log m,\log n)=x(\log2,\log3)$ does not define an element
$x\bmod\ell$ and does not imply residue compatibility.  In fact,
\cref{cor:generic-incompatibility} says that if a counterexample exists, then
an attempted common residue exponent fails at almost every Kummer prime as
$\ell$ grows.  Hence a proof cannot proceed by the sentence
``reduce the real exponent modulo $\ell$''; there is no such reduction.

A successful local proof would need one of the following genuinely new
bridges:
\begin{itemize}
\item a reason to select primes from the sparse compatible Chebotarev set;
\item a global product-formula argument that couples many incompatible
      places without assigning a local exponent;
\item or a theorem that turns the integrality of the entire orbit, rather
      than one pair of values, into a finite-place constraint.
\end{itemize}
The next section tests the third possibility exactly.

\section{The normalized mantissa orbit}\label{sec:mantissa}

\subsection{A dense rational orbit on a compact arc}

For $j\ge0$, define
\begin{equation}\label{eq:mantissa-def}
 k_j=\lfloor jx\rfloor,\qquad \theta_j=\{jx\},\qquad
 U_j=\frac{m^j}{2^{k_j}}=2^{\theta_j},\qquad
 V_j=\frac{n^j}{3^{k_j}}=3^{\theta_j}.
\end{equation}
Thus
\[
 1\le U_j<2,\qquad 1\le V_j<3,
\]
and every $(U_j,V_j)$ is a rational point on the compact analytic arc
\[
 \mathcal A=\{(2^t,3^t):0\le t<1\}.
\]
Since $x$ is irrational, $(\theta_j)$ is dense and uniformly distributed in
$[0,1)$; hence the rational points \eqref{eq:mantissa-def} are dense on
$\mathcal A$.

This reformulation packages all positive points of the additive lattice
$\Z+\Z x$.  It is stronger than looking at $(m,n)$ alone, but its height
also grows exponentially with $j$, so density by itself does not violate
rational-point counting theorems.

\subsection{Joint equidistribution of the index and the carry}

The pair $(j,k_j)$ is the discrete coding of the irrational rotation by
$x$.  The following exact refinement is the key input.

\begin{theorem}[conditional carry equidistribution]\label{thm:carry-UD}
Let $x$ be irrational and let $q\ge2$ be an integer.  For every residue pair
$(r,s)\in(\Z/q\Z)^2$ and every interval $I\subset[0,1)$,
\begin{equation}\label{eq:carry-density}
 \lim_{N\to\infty}\frac1N
 \#\{0\le j<N:
      j\equiv r\pmod q,
      \lfloor jx\rfloor\equiv s\pmod q,
      \{jx\}\in I\}
 =\frac{|I|}{q^2}.
\end{equation}
The same holds for finite unions of intervals whose boundaries have measure
zero.  Equivalently,
\[
 \bigl(\{jx\},j\bmod q,\lfloor jx\rfloor\bmod q\bigr)
\]
is equidistributed for the product of Lebesgue measure and uniform counting
measure on $(\Z/q\Z)^2$.
\end{theorem}

\begin{proof}
Fix $r$ and write $j=r+q h$.  Put
\[
 y_h=\frac{jx}{q}=\frac{rx}{q}+hx,
 \qquad t_h=\{y_h\}.
\]
Because $x$ is irrational, $(t_h)$ is uniformly distributed in $[0,1)$.
Now
\[
 jx=q\lfloor y_h\rfloor+q t_h,
\]
so
\[
 \lfloor jx\rfloor\bmod q=\lfloor q t_h\rfloor,
 \qquad \{jx\}=\{q t_h\}.
\]
The conditions $\lfloor jx\rfloor\equiv s\pmod q$ and $\{jx\}\in I$ are
therefore equivalent to
\[
 t_h\in\frac{s+I}{q}
   =\left\{\frac{s+u}{q}:u\in I\right\},
\]
an interval of length $|I|/q$.  Among indices with $j\equiv r\pmod q$ this
has conditional density $|I|/q$; the residue class $r$ itself has density
$1/q$.  Their product is $|I|/q^2$.
\end{proof}

The theorem is elementary, but it supplies the missing quantitative statement
that ``real closeness'' may be imposed simultaneously with arbitrary residues
of both $j$ and its carry $\lfloor jx\rfloor$.

\subsection{Residue symbols of the mantissa orbit}

Fix an odd prime $\ell$ and an unramified prime ideal $\mathfrak p$ as in
\cref{sec:kummer-setup}.  Multiplicativity of the Kummer character and
\eqref{eq:mantissa-def} give
\begin{equation}\label{eq:mantissa-kappa}
 \begin{pmatrix}
  \kappa_{\mathfrak p}(U_j)\\
  \kappa_{\mathfrak p}(V_j)
 \end{pmatrix}
 =
 M_{\mathfrak p}
 \begin{pmatrix}j\\k_j\end{pmatrix}
 \pmod\ell,
 \qquad
 M_{\mathfrak p}=
 \begin{pmatrix}
  \kappa_{\mathfrak p}(m)&-\kappa_{\mathfrak p}(2)\\
  \kappa_{\mathfrak p}(n)&-\kappa_{\mathfrak p}(3)
 \end{pmatrix}.
\end{equation}
Its determinant is
\begin{equation}\label{eq:mantissa-det}
 \det M_{\mathfrak p}
 =\kappa_{\mathfrak p}(2)\kappa_{\mathfrak p}(n)
  -\kappa_{\mathfrak p}(3)\kappa_{\mathfrak p}(m)
 =\Delta(\boldsymbol\kappa_{\mathfrak p}).
\end{equation}

\begin{theorem}[mantissa-orbit entropy]\label{thm:mantissa-entropy}
Let $(x,m,n)$ be a counterexample datum, let $\ell$ be an odd prime, and let
$\mathfrak p$ be an unramified prime ideal as above.  For every interval
$I\subset[0,1)$, conditional on $\theta_j\in I$ the vectors
\[
 \bigl(\kappa_{\mathfrak p}(U_j),
       \kappa_{\mathfrak p}(V_j)\bigr)
\]
are uniformly distributed on $\operatorname{im}M_{\mathfrak p}$.  More
precisely, each $w\in\operatorname{im}M_{\mathfrak p}$ occurs with natural
density
\[
 \frac{|I|}{|\operatorname{im}M_{\mathfrak p}|}.
\]
Consequently:
\begin{enumerate}
\item if $\Delta(\boldsymbol\kappa_{\mathfrak p})\ne0$, every one of the
      $\ell^2$ symbol pairs occurs with density $|I|/\ell^2$;
\item if $M_{\mathfrak p}$ has rank one, every point on its image line occurs
      with density $|I|/\ell$;
\item if $M_{\mathfrak p}=0$, the symbol pair is identically zero.
\end{enumerate}
\end{theorem}

\begin{proof}
By \cref{thm:carry-UD}, conditional on $\theta_j\in I$, the vector
$(j,k_j)\bmod\ell$ is uniform on $\F_\ell^2$.  A linear map sends uniform
measure to uniform measure on its image, with every image point having
$|\ker M_{\mathfrak p}|$ preimages.  Equation
\eqref{eq:mantissa-kappa} completes the proof.
\end{proof}

\begin{corollary}[full local entropy at generic primes]\label{cor:full-entropy}
For a counterexample datum and every sufficiently large good odd prime
$\ell$, the following holds.
\begin{enumerate}
\item In multiplicative rank four, a relative Chebotarev proportion
\[
 1-\frac1\ell-\frac1{\ell^2}+\frac1{\ell^3}
\]
of prime ideals have full mantissa-orbit entropy $\ell^2$ on every real
subinterval.
\item In multiplicative rank three, a relative proportion $1-1/\ell$ have
full entropy.
\end{enumerate}
\end{corollary}

\begin{proof}
Full entropy is equivalent to $\Delta\ne0$ by
\eqref{eq:mantissa-det}.  Apply \cref{prop:projective-law}.
\end{proof}

\subsection{Interpretation: a rigorous anti-continuity result}

The points $(U_j,V_j)$ can be restricted to an arbitrarily short real arc
around any point $(2^t,3^t)$, yet at a Chebotarev-generic prime their residue
symbols still cover all of $\mu_\ell^2$ uniformly.  Thus no argument of the
following form can be valid:
\begin{quote}
The real points lie nearly on one tangent direction, so their reductions
must have approximately one-dimensional residue data.
\end{quote}
There is no finite-field notion of ``approximately,'' and the exact theorem
shows two-dimensional behavior even after arbitrarily severe archimedean
conditioning.

This does not make local methods useless.  It says that any useful local
constraint must involve heights, divisibility by primes that vary with $j$,
or a simultaneous adelic product---not reduction modulo one fixed prime.
The selected-prime strategy in \cref{sec:selected-prime} follows precisely
that prescription.

\section{A rigidity barrier for auxiliary polynomials}\label{sec:aux-rigidity}

Auxiliary-polynomial and interpolation-determinant methods are natural here.
A counterexample supplies algebraic points
\[
 (2^{a+bx},3^{a+bx})=(2^a m^b,3^a n^b)
\]
for all integers $a,b$ for which the displayed quantities are interpreted as
rationals, and as integers when $a,b\ge0$.  One might hope to use the roughly
$D^2$ coefficients of a bivariate degree-$D$ polynomial to obtain roughly
$D^2$ orders of vanishing along the one-dimensional exponential curve.  The
next theorem shows that algebraic coefficients forbid exactly that gain.

\subsection{Exact equality of directional order and ordinary multiplicity}

For a polynomial $P\in\overline\Q[X,Y]$ and a point
$(A,B)\in\overline\Q^{\times2}$, let
$\mult_{(A,B)}P$ be the least total degree of a nonzero homogeneous term in
the Taylor expansion of $P(A+X,B+Y)$.  It is zero when $P(A,B)\ne0$.

\begin{theorem}[transcendental-direction rigidity]\label{thm:direction-rigidity}
Let $z_0\in\C$ be such that
\[
 A=2^{z_0},\qquad B=3^{z_0}
\]
are algebraic and nonzero.  For every nonzero
$P\in\overline\Q[X,Y]$, define
\[
 F(t)=P\bigl(Ae^{t\log2},Be^{t\log3}\bigr)
      =P(2^{z_0+t},3^{z_0+t}).
\]
Then
\begin{equation}\label{eq:ord-equals-mult}
 \ord_{t=0}F=\mult_{(A,B)}P.
\end{equation}
\end{theorem}

\begin{proof}
Let $d=\mult_{(A,B)}P$ and write
\[
 P(A+X,B+Y)=H_d(X,Y)+\text{terms of total degree greater than }d,
\]
where $H_d\in\overline\Q[X,Y]$ is a nonzero homogeneous polynomial of degree
$d$.  As $t\to0$,
\[
 Ae^{t\log2}-A=A(\log2)t+O(t^2),\qquad
 Be^{t\log3}-B=B(\log3)t+O(t^2).
\]
Therefore
\begin{align*}
 F(t)
 &=H_d\bigl(A\log2,B\log3\bigr)t^d+O(t^{d+1})\\
 &=(\log2)^d H_d(A,B\alpha)t^d+O(t^{d+1}),
 \qquad \alpha=\frac{\log3}{\log2}.
\end{align*}
The polynomial $T\mapsto H_d(A,BT)$ is a nonzero polynomial over
$\overline\Q$: multiplication of its coefficients by the nonzero algebraic
numbers $A$ and $B$ cannot annihilate it.  By
\cref{prop:alpha-transcendental}, it does not vanish at $T=\alpha$.  Hence the
coefficient of $t^d$ is nonzero, proving \eqref{eq:ord-equals-mult}.
\end{proof}

\begin{corollary}[directional conditions are all mixed conditions]\label{cor:mixed-conditions}
Under the hypotheses of \cref{thm:direction-rigidity}, the following are
equivalent for an integer $T\ge1$:
\begin{enumerate}
\item $F^{(j)}(0)=0$ for $0\le j<T$;
\item every mixed partial derivative
$\partial_X^r\partial_Y^sP(A,B)$ with $r+s<T$ vanishes;
\item $P$ has ordinary bivariate multiplicity at least $T$ at $(A,B)$.
\end{enumerate}
\end{corollary}

\begin{proof}
The equivalence of the second and third statements is the definition of
multiplicity in characteristic zero.  The first is $\ord_0F\ge T$, which is
equivalent to the third by \cref{thm:direction-rigidity}.
\end{proof}

\begin{corollary}[degree barrier]\label{cor:degree-barrier}
If $P$ has total degree at most $D$, then
\[
 \ord_{t=0}P(2^{z_0+t},3^{z_0+t})\le D.
\]
If $P$ has bidegree at most $(D,D)$, the order is at most $2D$.
\end{corollary}

\begin{proof}
The multiplicity of a nonzero polynomial at a point cannot exceed its total
degree.  A polynomial of bidegree $(D,D)$ has total degree at most $2D$.
\end{proof}

The theorem applies at every arithmetic lattice point $z_0=a+bx$, because
$2^{z_0}=2^am^b$ and $3^{z_0}=3^an^b$ are rational.  Hence the obstruction is
not confined to the base point $(1,1)$.

\subsection{Why the tempting dimension count fails}

The vector space of polynomials of total degree at most $D$ has dimension
$\binom{D+2}{2}\asymp D^2/2$.  If one viewed
$P(2^t,3^t)$ as an arbitrary one-variable analytic function, it would be
tempting to impose almost $D^2/2$ vanishing derivatives at one point.  But
those derivative equations have coefficients involving powers of
$\alpha=\log3/\log2$.  For algebraic coefficients of $P$, the transcendence
of $\alpha$ separates every coefficient of every power of $\alpha$.
Consequently the first $T$ directional equations are exactly the
$\binom{T+1}{2}$ mixed Taylor equations.  The apparent one-dimensional
problem retains its full two-dimensional arithmetic cost.

This explains a familiar threshold without appealing to vague statements
about ``not enough parameters.''  A successful auxiliary construction must
do something genuinely different from exact high contact by an algebraic
polynomial.  Possibilities include:
\begin{itemize}
\item approximate rather than exact contact, together with an unusually
      strong algebraic-independence measure for $\alpha$;
\item simultaneous smallness at many points with denominators that cancel in
      a global determinant;
\item an auxiliary object outside $\overline\Q[X,Y]$, but with enough
      arithmetic structure to retain a nonzero lower bound;
\item or a finite-place divisibility gain that replaces high archimedean
      multiplicity.
\end{itemize}

\subsection{A Wronskian companion}

For completeness, let $(i_1,j_1),\ldots,(i_N,j_N)$ be distinct pairs of
nonnegative integers and put
\[
 f_\nu(z)=2^{i_\nu z}3^{j_\nu z}.
\]
Their Wronskian is
\begin{equation}\label{eq:wronskian}
 W(f_1,\ldots,f_N)(z)
 =\left(\prod_{\nu=1}^N f_\nu(z)\right)
   \prod_{1\le\mu<\nu\le N}
   \bigl((i_\nu-i_\mu)\log2+(j_\nu-j_\mu)\log3\bigr),
\end{equation}
up to the sign determined by the ordering.  It never vanishes, since a zero
factor would make $\alpha$ rational.  Thus these monomials form an extended
Chebyshev system on the real line: a nontrivial linear combination with $N$
terms has at most $N-1$ real zeros counted with multiplicity.  This analytic
zero estimate is sharp over arbitrary complex coefficients, while
\cref{thm:direction-rigidity} is the stronger arithmetic local statement for
algebraic polynomial coefficients.

\begin{warning}[the auxiliary-polynomial trap]\label{warn:aux-trap}
A claim of order $\gg D$ for
$P(2^{z_0+t},3^{z_0+t})$ with
$P\in\overline\Q[X,Y]$ of degree $D$ is false unless $P=0$.  If the proposed
construction uses $D^2$ coefficients to impose $D^2$ directional derivatives,
then either the coefficient field secretly contains $\alpha$, the equations
are only approximate, or the linear system has only the zero solution.
\end{warning}

\section{Near collisions and a selected-prime bridge}\label{sec:selected-prime}

The Kummer theorem says generic fixed primes are hostile to a local exponent.
A more plausible strategy is to let the prime vary with a rational
approximation to $x$.  This section identifies exactly where such primes are
forced to lie.

\subsection{Simultaneous near-power differences}

Let $a,b\in\Z_{>0}$ and put
\[
 \delta=bx-a.
\]
Then
\begin{equation}\label{eq:near-differences}
 E_{a,b}=m^b-2^a=2^a(2^\delta-1),\qquad
 F_{a,b}=n^b-3^a=3^a(3^\delta-1).
\end{equation}
When $a/b$ is a continued-fraction convergent to $x$, one has
$|\delta|\ll1/b$.  For $|\delta|\le1$, the mean-value theorem gives constants
$c_1,c_2>0$, independent of $a,b$, such that
\begin{equation}\label{eq:near-bounds}
 c_1 2^a|\delta|\le |E_{a,b}|\le c_2 2^a|\delta|,
 \qquad
 c_1 3^a|\delta|\le |F_{a,b}|\le c_2 3^a|\delta|.
\end{equation}
Thus the differences are small relatively, but still exponentially large
absolutely.  Any argument that concludes $E_{a,b}=0$ merely from
$|\delta|\ll1/b$ has lost an exponential factor.

\subsection{Common divisors are Kummer-compatible}

\begin{proposition}[common-divisor localization]\label{prop:gcd-localization}
Let $\ell$ be an odd prime.  Let $p\nmid6mn\ell$ be a rational prime with
$p\equiv1\pmod\ell$, and let $\mathfrak p$ be a prime of $K_\ell$ above $p$.
Suppose
\[
 p\mid E_{a,b},\qquad p\mid F_{a,b},\qquad \ell\nmid b.
\]
Then $\mathfrak p$ is exactly compatible, with common residue exponent
\[
 t\equiv a b^{-1}\pmod\ell.
\]
\end{proposition}

\begin{proof}
The divisibilities give
\[
 m^b\equiv2^a\pmod{\mathfrak p},\qquad
 n^b\equiv3^a\pmod{\mathfrak p}.
\]
Applying the order-$\ell$ residue character yields
\[
 b\kappa_{\mathfrak p}(m)=a\kappa_{\mathfrak p}(2),\qquad
 b\kappa_{\mathfrak p}(n)=a\kappa_{\mathfrak p}(3)
 \quad\text{in }\F_\ell.
\]
Since $b$ is invertible modulo $\ell$, these are exactly the two compatibility
equations with $t=ab^{-1}$.
\end{proof}

\begin{corollary}[Chebotarev sparsity of common factors]\label{cor:gcd-sparse}
For a counterexample datum and a sufficiently large good $\ell$, every prime
$p\equiv1\pmod\ell$ satisfying the hypotheses of
\cref{prop:gcd-localization} lies over the exact compatible Chebotarev set.
Its relative density is
\[
 \frac1\ell-\frac1{\ell^2}+\frac1{\ell^3}
\]
in rank three and
\[
 \frac1\ell-\frac1{\ell^3}+\frac1{\ell^4}
\]
in rank four.
\end{corollary}

This is the cleanest connection found between archimedean approximation and
the Kummer detector.  It does not assign a residue exponent to $x$; rather,
a prime that divides both approximation errors manufactures the exponent
$t=a/b$ modulo $\ell$.

\subsection{Why density alone is insufficient}

Chebotarev describes primes in the ambient set, not prime divisors of one
thin sequence.  A set of density $1/\ell$ can contain every prime divisor of
all the integers $\gcd(E_{a,b},F_{a,b})$ without contradiction.  Neither the
prime number theorem nor effective Chebotarev supplies a ``generic prime
factor'' of a prescribed integer.

Likewise, the $abc$ conjecture applied separately to
\[
 m^b=2^a+E_{a,b},\qquad n^b=3^a+F_{a,b}
\]
would force substantial radical in each error term, but would not force the
two radicals to overlap or to escape a specified Frobenius set.  Known upper
bounds for greatest common divisors of fixed-base sequences, such as the
Bugeaud--Corvaja--Zannier theorem, point in the opposite direction: they make
large common divisors exceptional.  Here a proof needs either a lower bound
or a prime-factor distribution theorem for two correlated, nonstationary
exponential differences.

\subsection{A precise sparse-support target}

The following target is intentionally stated at the strength actually needed,
not as an informal appeal to ``random prime factors.''

\begin{target}[Frobenius escape for simultaneous near powers]\label{target:frob-escape}
Let $m,n\in\Z_{>1}$ and let
\[
 x=\frac{\log m}{\log2}=\frac{\log n}{\log3}\notin\Q.
\]
Then there exists a sufficiently large odd prime $\ell$ and infinitely many
continued-fraction convergents $a/b$ to $x$, with $\ell\nmid b$, such that
$\gcd(E_{a,b},F_{a,b})$ has a prime divisor $p\equiv1\pmod\ell$ whose
Frobenius vector is not in $\mathcal C_\ell$.
\end{target}

\begin{proposition}\label{prop:escape-finishes}
\cref{target:frob-escape} implies the Alaoglu--Erd\H{o}s conjecture.
\end{proposition}

\begin{proof}
Under a counterexample datum, the prime supplied by the target would satisfy
the hypotheses of \cref{prop:gcd-localization}, which says its Frobenius
vector is in $\mathcal C_\ell$.  This contradicts the conclusion of the
target.
\end{proof}

The target may be stronger than necessary, because the gcd can often be one.
A more flexible version would replace an actual common divisor by an adelic
linear combination of $E_{a,b}$ and $F_{a,b}$ whose product formula forces
Frobenius escape.  No such construction is presently available.  Still,
\cref{target:frob-escape} gives a concrete interface for analytic number
theory: it asks for a prime divisor in a named conjugacy subset of an explicit
pair of integers.

\section{The one-more-base route}\label{sec:third-base}

A different way to break the $2\times2$ barrier is to manufacture one more
multiplicatively independent algebraic base whose $x$th power is algebraic.
The Six Exponentials Theorem then finishes the problem immediately.

\begin{definition}
For a fixed real $x$, define the algebraic-power domain
\[
 \mathcal D_x=
 \{a\in\overline\Q_{>0}^{\times}:a^x\in\overline\Q^{\times}\},
\]
using the positive real value when $a>0$.
\end{definition}

\begin{theorem}[third-base exclusion]\label{thm:third-base-exclusion}
Let $(x,m,n)$ be a counterexample datum.  If
$a\in\overline\Q_{>0}^{\times}$ and $a^x$ is algebraic, then the three
logarithms $\log2,\log3,\log a$ are linearly dependent over $\Q$.
In particular,
\begin{equation}\label{eq:Dx-integers}
 \mathcal D_x\cap\Z_{>0}
 \subseteq\{2^r3^s:r,s\in\Z_{\ge0}\}.
\end{equation}
\end{theorem}

\begin{proof}
If the three logarithms were $\Q$-linearly independent, use them as the
three first factors in the Six Exponentials Theorem and use the two
$\Q$-linearly independent numbers $1,x$ as the second factors.  The six
exponentials would be
\[
 2,\ 3,\ a,\ 2^x=m,\ 3^x=n,\ a^x,
\]
all algebraic, contradicting that theorem.  For an integer $a$, a
multiplicative dependence among $2,3,a$ means that no prime outside
$\{2,3\}$ divides $a$, which proves \eqref{eq:Dx-integers}.
\end{proof}

\begin{corollary}[failure of immediate iteration]\label{cor:iteration-fails}
For a counterexample datum, at least one of $m^x=2^{x^2}$ and
$n^x=3^{x^2}$ is transcendental.
\end{corollary}

\begin{proof}
By \cref{cor:outside-prime}, at least one of $m,n$ has a prime divisor outside
$\{2,3\}$.  If its $x$th power were algebraic, it would violate
\cref{thm:third-base-exclusion}.
\end{proof}

Thus the most obvious attempt to iterate the map $a\mapsto a^x$ is blocked in
at least one coordinate.  Conversely, this makes the proof target extremely
sharp.

\begin{target}[one-more-base lemma]\label{target:one-more-base}
Assume $2^x,3^x\in\Z$ for a real $x$.  Construct an integer
$a$ with a prime divisor outside $\{2,3\}$ such that $a^x$ is algebraic.
\end{target}

\begin{proposition}
\cref{target:one-more-base} implies \cref{conj:AE}.
\end{proposition}

\begin{proof}
If $x$ were nonintegral, it would be irrational and the constructed $a$
would contradict \cref{thm:third-base-exclusion}.
\end{proof}

The difficulty is that exponentiation respects multiplication but not
addition, factor extraction, or polynomial operations.  From
$m=2^x,n=3^x$ one obtains $(2^r3^s)^x=m^rn^s$, but these bases remain in the
rank-two group generated by $2$ and $3$.  A valid third-base argument must
supply a new multiplicative direction without using a false identity such as
$(m+n)^x=m^x+n^x$.

\section{Three proof bridges and one no-go theorem}\label{sec:bridges}

The preceding sections turn several broad ideas into exact interfaces.  This
section collects them in a form suitable for a research program or a proof
assistant backlog.

\subsection{Bridge I: universal local compatibility}

The strongest possible local bridge would say that the real exponent relation
forces compatibility at every prime.  One good Kummer level would already be
enough.

\begin{target}[universal Kummer bridge]\label{target:universal-kummer}
For some odd prime $\ell$ outside the saturation exceptional set, prove from
\[
 \frac{\log m}{\log2}=\frac{\log n}{\log3}
\]
that every unramified Kummer Frobenius vector in $W_\ell$ belongs to
$\mathcal C_\ell$.
\end{target}

\begin{proposition}\label{prop:universal-finishes}
\cref{target:universal-kummer} implies the Alaoglu--Erd\H{o}s conjecture.
\end{proposition}

\begin{proof}
If a nonintegral counterexample had rank four, the compatibility density
would be the proper fraction in \eqref{eq:rank4-compatible}; if it had rank
three, it would be the proper fraction in \eqref{eq:rank3-compatible}.
Chebotarev guarantees Frobenius elements outside $\mathcal C_\ell$ in either
case, contradicting the target.  Equivalently, a linear subspace of
$2\times2$ matrices of dimension at least three cannot consist entirely of
singular matrices, while $\mathcal C_\ell$ is contained in the singular
cone.
\end{proof}

\cref{thm:mantissa-entropy} shows why a continuity proof of this target cannot
work: at most primes the normalized orbit has full residue entropy even on a
short real arc.  Any proof of the universal bridge would therefore have to
use a global arithmetic invariant not visible in fixed-prime reduction.

\subsection{Bridge II: selected-prime Frobenius escape}

This is \cref{target:frob-escape}.  It asks varying prime divisors of the
simultaneous near-power differences to escape a Chebotarev set of relative
density about $1/\ell$.  It is the most concrete interface with analytic
number theory, sieve theory, and arithmetic dynamics.

A potentially more attainable variant is the following product-formula
version.

\begin{target}[adelic escape]\label{target:adelic-escape}
For a counterexample datum, construct from infinitely many convergents $a/b$
a nonzero rational or algebraic number $G_{a,b}$ satisfying all three
properties:
\begin{enumerate}
\item its archimedean absolute value is bounded above by
      $\exp(-c\,h(G_{a,b}))$ for some fixed $c>0$;
\item at primes in the compatible Kummer set, its negative valuations have a
      total contribution $o(h(G_{a,b}))$;
\item at incompatible primes, \cref{prop:gcd-localization} or a replacement
      forces nonnegative valuation.
\end{enumerate}
Then the product formula gives a contradiction for large height.
\end{target}

This target makes explicit what a vague ``combine the real and $p$-adic
estimates'' must actually deliver.  The hard part is the second condition:
Chebotarev density by itself does not control the valuation mass of a thin
integer sequence.

\subsection{Bridge III: one more algebraic power}

This is \cref{target:one-more-base}.  It is algebraically the cleanest route:
one new multiplicative direction invokes Six Exponentials and ends the
problem.  Candidate constructions should be screened by the following
closure test.

\begin{proposition}[closure test]\label{prop:closure-test}
Let $G=\langle2,3\rangle\subset\Q_{>0}^{\times}$.  Every base obtained from
$2$ and $3$ using multiplication, inversion, and rational powers remains in
$G\otimes\Q$ and cannot serve as the third base.  Every base obtained from
$m,n$ by the same operations has a known $1/x$th power, not in general a
known $x$th power.  Therefore a successful construction must either make the
$x$-power of an output-side element algebraic or use an operation beyond the
multiplicative group while retaining exact control of exponentiation.
\end{proposition}

\begin{proof}
The first statement is closure of a divisible hull.  For the second, if
$a=m^rn^s$, then $a^{1/x}=2^r3^s$ is algebraic, whereas
$a^x=m^{rx}n^{sx}$ contains the unknown quantities $m^x,n^x$.  No
multiplicative manipulation removes them unless an additional relation is
proved.
\end{proof}

\subsection{No-go theorem: exact high directional contact}

Theorem~\ref{thm:direction-rigidity} closes a broad family of auxiliary
attacks.  It says that a target of the form
\[
 \ord_0 P(2^t,3^t)\gg\deg P
\]
with algebraic polynomial coefficients is not merely difficult; it is
false.  A revised auxiliary method must use approximate smallness, multiple
points, finite-place divisibility, or a nonpolynomial arithmetic object.

\subsection{Architecture scorecard}

\begin{longtable}{@{}p{0.19\textwidth}p{0.23\textwidth}p{0.25\textwidth}p{0.23\textwidth}@{}}
\toprule
Architecture & Rigorous gain in this report & Exact missing statement & Main failure mode to avoid\\
\midrule
\endfirsthead
\toprule
Architecture & Rigorous gain in this report & Exact missing statement & Main failure mode to avoid\\
\midrule
\endhead
Kummer / Chebotarev & Exact compatible and determinant densities; rank-three transversality & Universal Kummer bridge or a selected compatible-prime mechanism & Treating a real exponent as an element of $\F_\ell$ or $\Z_p$\\
\addlinespace
Mantissa dynamics & Conditional real--finite equidistribution and full local entropy & A height-sensitive restriction coupling many places & Inferring finite continuity from a short real arc\\
\addlinespace
Near-power gcd & Every common divisor is localized to the compatible set & Frobenius escape for prime divisors of correlated errors & Applying ambient prime density to a thin sequence\\
\addlinespace
Auxiliary polynomial & Exact equality of directional order and ordinary multiplicity & An approximate or adelic replacement with a nonzero lower bound & Counting $D^2$ coefficients as $D^2$ directional zeros\\
\addlinespace
Six Exponentials & Third-base exclusion and an immediate finishing criterion & Produce one new base $a$ with $a^x$ algebraic & Reusing a base in the divisible hull of $\langle2,3\rangle$\\
\addlinespace
Rank-three structure & Möbius normal form and transverse relation determinant & Rule out the remaining fractional-linear case & Assuming transcendence of $\alpha$ implies algebraic independence of arbitrary functions of $\alpha$\\
\bottomrule
\end{longtable}

\section{Finite verification and reproducible experiments}\label{sec:computation}

The density theorems are conceptual consequences of Kummer theory and
Chebotarev, but their finite combinatorics should be exhaustively checked
before formalization.  The following table gives the exact cardinalities for
small odd primes.

\begin{center}
\begin{tabular}{@{}rrrrr@{}}
\toprule
$\ell$ & $|\mathcal C_\ell|$ in rank 4 & transverse rank-3 compatible & rank-4 $\Delta=0$ & rank-3 $\Delta=0$\\
\midrule
3  & 25   & 7   & 33   & 9\\
5  & 121  & 21  & 145  & 25\\
7  & 337  & 43  & 385  & 49\\
11 & 1321 & 111 & 1441 & 121\\
\bottomrule
\end{tabular}
\end{center}
The ambient cardinalities are respectively $\ell^4$ for rank four and
$\ell^3$ for rank three.

\subsection{Exhaustive checker}

\begin{lstlisting}[language=Python,caption={Exhaustive verification of the compatibility formulas}]
from itertools import product


def compatible(q: int, z: tuple[int, int, int, int]) -> bool:
    u1, u2, v1, v2 = z
    return any(
        (v1 - t * u1) % q == 0 and (v2 - t * u2) % q == 0
        for t in range(q)
    )


def det_zero(q: int, z: tuple[int, int, int, int]) -> bool:
    u1, u2, v1, v2 = z
    return (u1 * v2 - u2 * v1) % q == 0


def in_hyperplane(q: int, rho, z) -> bool:
    return sum(a * b for a, b in zip(rho, z)) % q == 0


def verify(q: int) -> None:
    all4 = list(product(range(q), repeat=4))
    comp = [z for z in all4 if compatible(q, z)]
    assert len(comp) == q**3 - q + 1

    singular = [z for z in all4 if det_zero(q, z)]
    assert len(singular) == q**3 + q**2 - q

    # Check every nonzero rank-three relation vector rho.
    for rho in all4[1:]:
        W = [z for z in all4 if in_hyperplane(q, rho, z)]
        Wc = [z for z in W if compatible(q, z)]
        A, B, C, D = rho
        transverse = (A * D - B * C) % q != 0
        expected = q**2 - q + 1 if transverse else None
        if transverse:
            assert len(Wc) == expected
            assert sum(det_zero(q, z) for z in W) == q**2


for q in (3, 5, 7, 11):
    verify(q)
print("all finite checks passed")
\end{lstlisting}

The code deliberately defines exact compatibility first and determinant
vanishing second.  That ordering prevents the false zero-source equivalence
in \cref{warn:det-not-exact}.

\subsection{Arithmetic experiments suggested by the theory}

There is no counterexample to sample.  Nevertheless, one can probe the
interfaces on arbitrary multiplicatively ranked tuples $(2,3,m,n)$ and on
near-relations.  The most informative experiments are:

\begin{enumerate}
\item \textbf{Kummer histogram.}  For fixed $\ell$, factor rational primes
$p\equiv1\pmod\ell$, compute order-$\ell$ residue-symbol coordinates, and
compare compatibility and determinant frequencies with the formulas for the
observed multiplicative rank.
\item \textbf{Rank-three transversality stress test.}  Generate tuples with a
known primitive relation, separate the cases $AD-BC=0$ and $AD-BC\ne0$, and
observe the two hyperplane counts in \cref{prop:hyperplane-count}.  The
proportional case should always correspond to a rational scalar in the
archimedean model.
\item \textbf{Mantissa entropy.}  For an arbitrary irrational
$x=\log m/\log2$ (without requiring $3^x$ integral), compare the empirical
symbol distribution of $U_j$ with the exact linear image of $(j,k_j)$.  For
a synthetic second coordinate, the determinant controls the transition from
a line to the full plane.
\item \textbf{Near-power prime factors.}  For pairs $(m,n)$ making the two log
ratios unusually close, factor
$E_{a,b},F_{a,b}$ along convergents and record Frobenius classes of common
prime factors.  This does not prove escape, but it can falsify overly strong
versions of \cref{target:frob-escape} and identify realistic hypotheses.
\item \textbf{Auxiliary contact audit.}  Randomly generate algebraic
polynomials with prescribed multiplicity at an arithmetic orbit point and
verify symbolically that the restricted order equals that multiplicity.  Any
higher observed order signals a bug in expansion or an accidental zero
polynomial.
\end{enumerate}

For large Kummer computations, SageMath or PARI/GP is preferable to a custom
discrete-log implementation.  The finite count checker above is independent
of number-field software and should be included in continuous integration.

\section{Lean formalization plan}\label{sec:lean}

A productive formalization should not begin with Chebotarev.  The results
separate into layers with sharply different library requirements.

\subsection{Layer 0: finite combinatorics}

This layer needs only finite fields and cardinalities.

\begin{enumerate}
\item Define
\[
 \mathrm{Compatible}(u,v)\;:\!\Longleftrightarrow\;
 \exists t,\ v=t\,u
\]
for $u,v\in(\mathrm{ZMod}\ \ell)^2$.
\item Prove $|\mathcal C_\ell|=\ell^3-\ell+1$ by splitting on $u=0$.
\item Define fixed-$t$ planes and prove their pairwise intersection is the
zero vector.
\item Prove the transverse hyperplane count $\ell^2-\ell+1$ by showing
$L+tM$ is nonzero for every $t$.
\item Prove the singular-matrix counts separately and derive the warning that
$\Delta=0$ is not exact compatibility.
\end{enumerate}

A schematic definition, not intended as drop-in syntax, is:

\begin{lstlisting}[caption={Lean-shaped statement of exact compatibility}]
def Compatible (u v : (ZMod ell) x (ZMod ell)) : Prop :=
  exists t : ZMod ell, v.1 = t * u.1 /\ v.2 = t * u.2

 theorem card_compatible [Fact (Nat.Prime ell)] :
  Fintype.card {z // Compatible z.1 z.2} = ell^3 - ell + 1
\end{lstlisting}

\subsection{Layer 1: valuation lattices and rank}

Represent a finitely supported positive rational by its prime-valuation
vector.  Formalize:
\begin{itemize}
\item rational exponents giving integers are integral;
\item the multiplicative-rank floor in \cref{prop:rank-floor};
\item saturation of a finitely generated sublattice and the finite exceptional
index $I_\Gamma$;
\item the primitive relation in rank three and the determinant argument in
\cref{lem:rank3-transverse}.
\end{itemize}
The only transcendence input in this layer is that
$\log3/\log2$ is transcendental.  Until Gelfond--Schneider is available in the
library, package this as a clearly named imported theorem rather than an
opaque axiom local to the proof.

\subsection{Layer 2: irrational rotation and mantissa entropy}

Formalize \cref{thm:carry-UD} by reducing it to uniform distribution of
$h\mapsto\{c+hx\}$.  The proof uses only affine transformations of an
irrational rotation and interval measures.  Then formalize the linear push-
forward in \cref{thm:mantissa-entropy}.  This layer is independent of
algebraic number theory if $\kappa(2),\kappa(3),\kappa(m),\kappa(n)$ are
abstract elements of a finite field.

This separation is valuable: the full entropy theorem can be machine-checked
as a statement about an arbitrary $2\times2$ matrix over $\F_\ell$ before
residue symbols are introduced.

\subsection{Layer 3: Kummer theory and Chebotarev}

Required ingredients are:
\begin{itemize}
\item cyclotomic fields containing $\mu_\ell$;
\item the Kummer isomorphism for a finite-dimensional subgroup of
      $K^\times/K^{\times\ell}$;
\item Frobenius elements and power-residue symbols;
\item Chebotarev density for finite Galois extensions;
\item the elementary injection of rational Kummer classes into
      $\Q(\zeta_\ell)$ from \cref{lem:eventual-rank}.
\end{itemize}
If the library lacks a sufficiently general Chebotarev theorem, a sensible
intermediate artifact is a theorem parameterized by an abstract
``Frobenius-uniformity'' hypothesis.  The finite formulas and every later
logical implication can then be formalized immediately; replacing the
hypothesis by Chebotarev becomes an isolated library project.

\subsection{Layer 4: analytic local rigidity}

The proof of \cref{thm:direction-rigidity} requires multivariate Taylor
expansion, order of a holomorphic zero, and transcendence of $\alpha$ over
$\overline\Q$.  A purely formal-power-series version may be easier:
expand
\[
 e^{t\log2}-1,\qquad e^{t\log3}-1
\]
and identify the lowest homogeneous Taylor term of $P$.  Only its linear
terms are needed to prove the leading coefficient nonzero.  This avoids a
large analytic dependency footprint.

\subsection{Recommended theorem order}

\begin{center}
\begin{tikzpicture}[
 node distance=8mm and 13mm,
 box/.style={draw,rounded corners,align=center,inner sep=5pt,text width=35mm},
 arr/.style={-{Latex[length=2mm]},thick}
]
\node[box] (finite) {Finite compatibility\\and determinant counts};
\node[box,right=of finite] (rank) {Valuation lattice\\and rank-three transversality};
\node[box,below=of finite] (rotation) {Carry equidistribution\\over $\F_\ell^2$};
\node[box,right=of rotation] (kummer) {Kummer coordinates\\and Chebotarev};
\node[box,below=of rotation] (entropy) {Mantissa-orbit\\entropy theorem};
\node[box,right=of entropy] (density) {Exact density laws\\for counterexample data};
\node[box,below=of density] (bridges) {Formal bridge lemmas\\and contradiction interfaces};

\draw[arr] (finite) -- (density);
\draw[arr] (rank) -- (density);
\draw[arr] (kummer) -- (density);
\draw[arr] (rotation) -- (entropy);
\draw[arr] (kummer) -- (entropy);
\draw[arr] (entropy) -- (bridges);
\draw[arr] (density) -- (bridges);
\end{tikzpicture}
\end{center}

The finite and dynamical branches can proceed in parallel.  They provide
useful checked artifacts even before the deepest algebraic-number-theory
layer is complete.

\section{Research agenda after this report}\label{sec:agenda}

\subsection{Highest-priority questions}

\paragraph{1. Can common-divisor localization be upgraded to Frobenius
escape?}
The exact proposition is now available.  The next step is to search the
literature on prime divisors of pairs of exponential-polynomial sequences,
especially results that prescribe a conjugacy class rather than merely a
congruence class.  A useful theorem must be uniform along convergents with
both the base and exponent normalization present.

\paragraph{2. Is there an adelic determinant whose forced divisibility beats
its archimedean size?}
The auxiliary-polynomial rigidity theorem says not to seek high exact contact
at one point.  A determinant involving several near-collision points may
still work if obvious powers of $2,3,m,n$ can be divided out while leaving a
nonzero algebraic integer.  Every proposed normalization should be audited
place by place; dividing by a large real factor without preserving
integrality is the standard hidden error.

\paragraph{3. Can one output be made to re-enter $\mathcal D_x$?}
By \cref{cor:iteration-fails}, a counterexample forces at least one immediate
second iterate to be transcendental.  A recurrence, algebraic correspondence,
or trace/norm construction that makes either $m^x$ or $n^x$ algebraic would
finish in many cases.  The construction must exploit integrality in a way
not preserved by arbitrary algebraic outputs, or else it would solve the
full Four Exponentials problem at once.

\paragraph{4. Can the rank-three Möbius form be attacked separately?}
In rank three,
\[
 x=-\frac{A+B\alpha}{C+D\alpha},\qquad AD-BC\ne0.
\]
This is a narrower special-function problem than unrestricted rank four.
Possible tools include transcendence measures for fractional-linear
functions of logarithm ratios, mixed real/$p$-adic regulators, and
Hermite--Pad\'e approximation constrained by the integer relation
$2^A3^Bm^Cn^D=1$.  Merely citing the transcendence of $\alpha$ is
insufficient: a rational function of a transcendental number can certainly
have algebraic exponentials.

\subsection{Directions now deprioritized}

The following directions are not impossible in principle, but the results
above show that their naive versions cannot be the missing proof.

\begin{itemize}
\item \textbf{A common $p$-adic exponent at an arbitrary prime.}  Generic
Kummer primes are incompatible with probability tending to one.
\item \textbf{Finite continuity from the real curve.}  Conditional mantissa
entropy is full at generic primes on every real subarc.
\item \textbf{Order-$D^2$ contact from a degree-$D$ algebraic polynomial.}
The exact local zero theorem bounds contact by $O(D)$.
\item \textbf{Chebotarev density as a statement about prime factors.}
Ambient density does not imply escape for a thin sequence.
\item \textbf{A third base built only by multiplying known bases.}  This
never leaves the rank-two divisible hull.
\end{itemize}

\subsection{A focused next experiment}

The most informative computational project is a two-level Frobenius audit.
Choose many integer pairs $(m,n)$ for which
$\log m/\log2$ and $\log n/\log3$ agree to unusually high precision but are
certifiably unequal.  For convergents $a/b$ to their mean, factor the common
part of $m^b-2^a$ and $n^b-3^a$.  At several Kummer levels $\ell$, record:

\begin{enumerate}
\item whether common factors lie in the exact compatible set as predicted;
\item how much logarithmic mass is carried by each Frobenius class;
\item whether the mass resembles ambient Chebotarev density or is strongly
      biased;
\item and whether bias correlates with multiplicative rank three versus four.
\end{enumerate}

The experiment cannot establish the conjecture, but it directly tests the
only bridge in this report that converts real approximants into selected
finite primes.  It should also expose whether
\cref{target:frob-escape} is realistic or needs an adelic reformulation.

\section{Proof-status ledger and conclusion}\label{sec:status}

\begin{longtable}{@{}p{0.42\textwidth}p{0.14\textwidth}p{0.35\textwidth}@{}}
\toprule
Statement & Status & Dependencies / caveat\\
\midrule
\endfirsthead
\toprule
Statement & Status & Dependencies / caveat\\
\midrule
\endhead
Rational $x$ with $2^x\in\Z$ is integral & proved & Unique factorization\\
Nonintegral counterexample has transcendental $x$ & proved & Gelfond--Schneider\\
$\alpha=\log3/\log2$ is transcendental & proved & Gelfond--Schneider\\
Multiplicative rank is at least three & proved & Transcendence of $\alpha$\\
Rank-three relation is transverse & proved & Rational case plus primitive relation\\
Eventual equality of ordinary and Kummer rank & proved & Valuation lattices and cyclotomic ramification\\
Universal compatibility-density formula & proved & Kummer theory and Chebotarev\\
Rank-four exact density & proved & Finite counting\\
Rank-three exact density & proved & Transversality and finite counting\\
Projective determinant densities & proved & Finite matrix counts\\
Conditional carry equidistribution & proved & Irrational rotation\\
Mantissa-orbit entropy theorem & proved & Carry equidistribution and character multiplicativity\\
Directional order equals ordinary multiplicity & proved & Transcendence of $\alpha$ and Taylor expansion\\
Common-divisor localization & proved & Residue-character congruences\\
Third-base exclusion & proved & Six Exponentials Theorem\\
Universal Kummer bridge & open target & Real relation gives no known finite exponent\\
Frobenius escape for near powers & open target & Requires prime-factor distribution in a thin pair\\
One-more-base lemma & open target & Must create a new multiplicative direction\\
Alaoglu--Erd\H{o}s conjecture & unresolved here & Any one of the three bridges would finish it\\
\bottomrule
\end{longtable}

The central lesson is not that local methods fail, but that their correct
shape is now visible.  A hypothetical counterexample is maximally hostile to
a common exponent at generic finite places: exact compatibility has density
about $1/\ell$, and the normalized orbit has full two-dimensional residue
entropy at almost all Kummer primes.  This rules out an entire family of
unjustified archimedean-to-$p$-adic transfers.

At the same time, simultaneous near-power differences provide a legitimate
selected-prime bridge, and Six Exponentials turns any third algebraic-power
base into an immediate proof.  The auxiliary-polynomial theorem eliminates
a false parameter gain and redirects interpolation work toward approximate,
multipoint, or adelic constructions.  These are substantive reductions of
the search space even though the final bridge remains open.

The strongest concrete next move is therefore not another broad survey.  It
is to attack the Frobenius-escape interface for
\[
 \gcd(m^b-2^a,n^b-3^a)
\]
along convergents, while formalizing the finite counts and mantissa entropy in
parallel.  A successful result there would connect the one place where the
real exponent genuinely manufactures a finite residue exponent to the place
where Chebotarev says such exponents are rare.

\appendix

\section{A compact derivation of the rank-three compatibility formula}

Let $W=\ker\rho\subset\F_\ell^4$ with
$\rho=(A,B,C,D)$ and $AD-BC\ne0$.  For each $t$, the fixed compatibility
plane is
\[
 \mathcal C_t=\{(u,tu):u\in\F_\ell^2\}.
\]
The equation for $W\cap\mathcal C_t$ is
\[
 (A+tC)u_1+(B+tD)u_2=0.
\]
The coefficient pair cannot vanish, because that would imply
$(A,B)=-t(C,D)$ and hence $AD-BC=0$.  Thus the intersection is a line with
$\ell$ points.  Distinct $t$ give lines meeting only at the origin.  Hence
\[
 \left|W\cap\bigcup_t\mathcal C_t\right|
 =1+\ell(\ell-1)=\ell^2-\ell+1.
\]
Dividing by $|W|=\ell^3$ gives
$\ell^{-1}-\ell^{-2}+\ell^{-3}$.  This argument is short enough to be the
preferred Lean proof.

\section{Saturation index from a prime-exponent matrix}

Let $S=\{p_1,\ldots,p_s\}$ contain the prime support of $6mn$.  Form the
$s\times4$ integer matrix
\[
 E=\bigl(v_{p_i}(a_j)\bigr),
 \qquad (a_1,a_2,a_3,a_4)=(2,3,m,n).
\]
Its column lattice is $L\subset\Z^s$.  A Smith normal form
\[
 UEV=\operatorname{diag}(d_1,\ldots,d_r,0,\ldots,0),
 \qquad d_1\mid d_2\mid\cdots\mid d_r,
\]
computes both the multiplicative rank $r$ and the saturation index
\[
 I_\Gamma=d_1\cdots d_r.
\]
For every prime $\ell\nmid I_\Gamma$, the reduction of the column span modulo
$\ell$ has dimension $r$.  This supplies an effective list of the Kummer
levels excluded by \cref{lem:eventual-rank}.  In rank three, also exclude
prime divisors of $AD-BC$.

\section{Bibliographic notes}

The original conjectural context is the work of Alaoglu and Erd\H{o}s on
highly composite and related integers.  The transcendence baseline is the
Gelfond--Schneider theorem and the Six Exponentials Theorem.  Kummer theory,
cyclotomic fields, and Chebotarev are used only in their standard finite
abelian form.  Uniform distribution enters through the elementary irrational
rotation theorem.  The selected-prime discussion is informed by the modern
literature on gcds of exponential sequences, but no such theorem is invoked
in a proof above.

\begin{thebibliography}{99}

\bibitem{AlaogluErdos1944}
L.~Alaoglu and P.~Erd\H{o}s,
\emph{On highly composite and similar numbers},
Transactions of the American Mathematical Society \textbf{56} (1944),
448--469.

\bibitem{Baker1975}
A.~Baker,
\emph{Transcendental Number Theory},
Cambridge University Press, 1975.

\bibitem{BombieriGubler2006}
E.~Bombieri and W.~Gubler,
\emph{Heights in Diophantine Geometry},
Cambridge University Press, 2006.

\bibitem{BCZ2003}
Y.~Bugeaud, P.~Corvaja, and U.~Zannier,
\emph{An upper bound for the G.C.D. of $a^n-1$ and $b^n-1$},
Mathematische Zeitschrift \textbf{243} (2003), 79--84.

\bibitem{EvertseSchlickeweiSchmidt2002}
J.-H.~Evertse, H.~P.~Schlickewei, and W.~M.~Schmidt,
\emph{Linear equations in variables which lie in a multiplicative group},
Annals of Mathematics \textbf{155} (2002), 807--836.

\bibitem{Gelfond1960}
A.~O.~Gelfond,
\emph{Transcendental and Algebraic Numbers},
Dover Publications, 1960.

\bibitem{KuipersNiederreiter1974}
L.~Kuipers and H.~Niederreiter,
\emph{Uniform Distribution of Sequences},
Wiley, 1974.

\bibitem{Lang1966}
S.~Lang,
\emph{Introduction to Transcendental Numbers},
Addison--Wesley, 1966.

\bibitem{Neukirch1999}
J.~Neukirch,
\emph{Algebraic Number Theory},
Springer, 1999.

\bibitem{SerreTopics}
J.-P.~Serre,
\emph{Topics in Galois Theory}, second edition,
A K Peters, 2008.

\bibitem{Waldschmidt2000}
M.~Waldschmidt,
\emph{Diophantine Approximation on Linear Algebraic Groups},
Springer, 2000.

\bibitem{Washington1997}
L.~C.~Washington,
\emph{Introduction to Cyclotomic Fields}, second edition,
Springer, 1997.

\end{thebibliography}

\end{document}
